\centerline{\bf \Huge Регион ВсОШ}
\centerline{\bf \Huge 10 класс}

\begin{flushright}
        \scriptsize{}
\end{flushright}

\problem{1}
Докажите, что у числа $2024!+2025$ есть простой делитель, больший 2026.

\problem{2}
На доску по кругу записано 2025 положительных чисел. Для каждых двух соседних чисел $a$ и $b$ Паша написал в тетрадку красным цветом число $a^2+b^2$ и синим цветом число $2 a b$. Докажите, что какое-то синее число не меньше какого-то красного.

\problem{3}
Ваня последовательно выставляет слонов на шахматную доску. Первого слона можно выставить произвольно, а каждый следующий в момент выставления должен бить нечётное число ранее выставленных слонов. Какое наибольшее количество слонов может выставить Ваня? (Слоны бьют друг друга, если они стоят на одной диагонали и между ними нет других слонов.)

\problem{4}
Многочлен $P(x)$ с целыми коэффициентами таков, что $P(0)=1$. Пусть $c-$ натуральное число, большее 1. Зададим последовательность $\left\{x_n\right\}$ следующим образом: $x_0=0, x_{i+1}=P\left(x_i\right)$ для всех $i \geqslant 0$. Докажите, что существует бесконечно много натуральных $n$, для которых НОД $\left(x_n, n+c\right)=1$.

\problem{5}
На гипотенузе $A B$ прямоугольного треугольника $A B C$ выбраны точки $M$ и $N$ так, что $A M<A N$. Прямая, проходящая через $M$ и перпендикулярная $C N$, пересекает прямую $A C$ в точке $P$. Прямая, проходящая через $N$ и перпендикулярная $C M$, пересекает прямую $B C$ в точке $Q$. Докажите, что описанные окружности треугольников $A P M$ и $B N Q$ и прямая $P Q$ имеют общую точку.
